\chapter{PostgreSQL --- Características Avanzadas}

\section{Por qué PostgreSQL}

RDBMS open source más avanzado. Extensible mediante extensiones, soporta JSON como ciudadano de primera clase, herencia de tablas, tipos custom, y tiene MVCC robusto.

\begin{itemize}
  \item JSONB: JSON binario indexable (vs JSON texto sin índice)
  \item Extensiones: PostGIS (geo), pg\_trgm (trigrams), uuid-ossp, pgcrypto
  \item VACUUM: recupera espacio de filas muertas (MVCC genera versiones viejas)
  \item Roles: sistema de permisos basado en roles, no usuarios individuales
  \item Particionamiento nativo: RANGE, LIST, HASH
\end{itemize}

\section{JSONB}

Almacena JSON en formato binario. Soporta índices GIN, operadores \texttt{->}, \texttt{->>}, \texttt{@>}, \texttt{?}.

\begin{itemize}
  \item \texttt{->} devuelve JSON (objeto/array)
  \item \texttt{->>} devuelve texto
  \item \texttt{@>} contiene (compatible con índice GIN)
  \item \texttt{?} tiene la clave
\end{itemize}

\section{VACUUM y mantenimiento}

VACUUM recupera espacio de tuplas muertas. AUTOVACUUM lo hace automáticamente. VACUUM ANALYZE actualiza estadísticas del planificador. VACUUM FULL bloquea la tabla (evitar en producción).

\begin{lstlisting}[caption={JSONB --- consultas y operadores}]
CREATE TABLE eventos (
    id      SERIAL PRIMARY KEY,
    datos   JSONB NOT NULL,
    creado  TIMESTAMP DEFAULT NOW()
);

INSERT INTO eventos (datos) VALUES
  ('{"tipo": "login", "usuario": "ana", "ip": "192.168.1.1"}'),
  ('{"tipo": "compra", "usuario": "carlos", "monto": 350.00}');

-- Acceder a campos
SELECT datos->>'usuario' AS usuario, datos->>'tipo' AS tipo
FROM eventos;

-- Filtrar con @> (usa índice GIN)
CREATE INDEX idx_eventos_gin ON eventos USING GIN(datos);
SELECT * FROM eventos WHERE datos @> '{"tipo": "login"}';

-- Actualizar campo en JSONB
UPDATE eventos
SET datos = datos || '{"procesado": true}'::jsonb
WHERE datos->>'tipo' = 'compra';
\end{lstlisting}

\begin{lstlisting}[caption={Roles, permisos y VACUUM}]
-- Crear rol de solo lectura
CREATE ROLE readonly;
GRANT CONNECT ON DATABASE mibd TO readonly;
GRANT USAGE ON SCHEMA public TO readonly;
GRANT SELECT ON ALL TABLES IN SCHEMA public TO readonly;

-- Crear usuario con ese rol
CREATE USER reportes_user WITH PASSWORD 'segura123';
GRANT readonly TO reportes_user;

-- VACUUM manual
VACUUM ANALYZE empleados;

-- Ver tablas que necesitan vacuum
SELECT relname, n_dead_tup, last_vacuum
FROM   pg_stat_user_tables
WHERE  n_dead_tup > 10000
ORDER BY n_dead_tup DESC;
\end{lstlisting}

\begin{lstlisting}[caption={Particionamiento nativo en PostgreSQL}]
-- Tabla particionada por rango de fecha
CREATE TABLE ventas (
    id      SERIAL,
    fecha   DATE NOT NULL,
    monto   NUMERIC(10,2),
    region  VARCHAR(50),
    PRIMARY KEY (id, fecha)
) PARTITION BY RANGE (fecha);

-- Crear particiones
CREATE TABLE ventas_2023 PARTITION OF ventas
    FOR VALUES FROM ('2023-01-01') TO ('2024-01-01');

CREATE TABLE ventas_2024 PARTITION OF ventas
    FOR VALUES FROM ('2024-01-01') TO ('2025-01-01');

-- PostgreSQL enruta automáticamente al insertar
INSERT INTO ventas (fecha, monto, region)
VALUES ('2024-06-15', 5000.00, 'Norte');
\end{lstlisting}

\begin{entrevistabox}
\begin{enumerate}
  \item ¿Cuál es la diferencia entre JSON y JSONB en PostgreSQL?
  \item ¿Qué hace VACUUM y por qué es necesario en PostgreSQL?
  \item ¿Cómo implementarías seguridad de solo lectura para un usuario de reportes?
\end{enumerate}
\end{entrevistabox}
